\section*{Abstract}

Die T0-Theorie postuliert weder den Parameter $\xi$ noch die minimale Längenskala $L_0 = \xi \cdot \lP$. Beide emergieren als notwendige Konsequenzen aus der Selbstkonsistenz des T0-Lagrangians. Dieses Dokument zeigt die vollständige Ableitungskette in fünf Schritten: (1) Lagrangian-Struktur mit Zeit-Masse-Dualität, (2) Selbstkonsistenzbedingung und $10^{-4}$-Skalierung, (3) Eindeutigkeit von $\xi$ aus dem e-p-$\mu$-System, (4) Gravitationskonstante $G$ aus $\xi$, (5) minimale Längenskala $L_0 = \xi \cdot \lP$.

\section{Schritt 1: Der T0-Lagrangian}

\subsection{Standardmodell-Ausgangspunkt}

Der Standardmodell-Lagrangian für ein Lepton $\psi_\ell$ lautet:
\begin{equation}
	\Lag_{\text{SM}} = -\frac{1}{4} F_{\mu\nu}F^{\mu\nu} + \bar{\psi}_\ell(i\gamma^\mu D_\mu - m_\ell)\psi_\ell
	\label{eq:sm_lagrangian}
\end{equation}
mit $D_\mu = \partial_\mu + ieA_\mu$ und konstantem $m_\ell$.

\subsection{Zeit-Masse-Dualität als fundamentales Prinzip}

In der T0-Theorie gilt die fundamentale Relation:
\begin{equation}
	\tilde{T}(x,t) \cdot m(x,t) = 1
	\label{eq:duality}
\end{equation}
Diese Relation ist kein zusätzliches Postulat, sondern eine strukturelle Bedingung: jede Masse trägt ein intrinsisches Zeitfeld. Daraus folgt zwingend ein dynamisches Massenfeld:
\begin{equation}
	m(x,t) = m_0 + \Delta m(x,t)
\end{equation}

\subsection{Vollständiger T0-Lagrangian}

\begin{keyresult}[Vollständiger T0-Lagrangian]
\begin{equation}
	\Lag_{\text{T0}} = -\frac{1}{4} F_{\mu\nu}F^{\mu\nu}
	+ \bar{\psi}_\ell(i\gamma^\mu D_\mu - m_\ell)\psi_\ell
	+ \frac{1}{2}(\partial_\mu \Delta m)(\partial^\mu \Delta m)
	- \frac{1}{2} m_T^2 \Delta m^2
	+ \xipar \, m_\ell \,\bar{\psi}_\ell \psi_\ell \, \Delta m
	\label{eq:t0_lagrangian}
\end{equation}
\end{keyresult}

\subsection{Die drei Schlüsselformeln}

\textbf{Kopplungsstärke} (massenproportional):
\begin{equation}
	\boxed{g_T^\ell = \xipar \cdot m_\ell}
	\label{eq:coupling}
\end{equation}

\textbf{Zeitfeldmasse:}
\begin{equation}
	m_T = \frac{\lambda}{\xipar}
\end{equation}

\textbf{Vertex-Faktor:}
\begin{equation}
	\bar{\psi}_\ell \psi_\ell \Delta m \longrightarrow -i\,\xipar\, m_\ell
\end{equation}

\textbf{Zeitfeld-Propagator:}
\begin{equation}
	\Delta m(k) \longrightarrow \frac{i}{k^2 - m_T^2 + i\epsilon}
\end{equation}

\section{Schritt 2: Selbstkonsistenzbedingung}

\subsection{Warum der Lagrangian \texorpdfstring{$\xi$}{xi} erzwingt}

Der Kopplungsterm $g_T^\ell = \xipar \cdot m_\ell$ ist nicht frei wählbar. Die Selbstkonsistenz des Lagrangians verlangt:
\begin{itemize}
	\item Die Theorie bleibt divergenzfrei bei hohen Energien.
	\item Die bekannten Teilchenmassen werden reproduziert.
	\item Die Massenhierarchie zwischen den Generationen ist geometrisch beschreibbar.
\end{itemize}

\subsection{Der \texorpdfstring{$10^{-4}$}{10^(-4)}-Faktor aus der QFT-Schleifenstruktur}

Die Schleifensuppression in vier Raumzeit-Dimensionen liefert:
\begin{equation}
	\frac{1}{16\pi^3} \approx 2{,}01 \times 10^{-3}
\end{equation}
Kombiniert mit den Higgs-Parametern aus dem T0-Lagrangian:
\begin{equation}
	(\lambda_h)^2 \cdot \frac{v^2}{m_h^2} = (0{,}129)^2 \times \frac{(246{,}2)^2}{(125{,}1)^2} \approx 0{,}0647
\end{equation}

\begin{keyresult}[\texorpdfstring{$10^{-4}$}{10^(-4)} aus der Raumzeit-Dimensionalität]
\begin{equation}
	\xipar_4 \sim (10^{-1})^4 = 10^{-4}
\end{equation}
Dieser Faktor ist nicht postuliert — er folgt aus der Dimensionalität der Raumzeit.
\end{keyresult}

\subsection{Der 4/3-Faktor aus der Tetraeder-Geometrie}

Der Faktor $4/3$ emergiert aus der harmonischen Struktur des dreidimensionalen Raums:
\begin{equation}
	\frac{4 \text{ Ecken}}{3 \text{ Kanten pro Ecke}} = \frac{4}{3} \quad \text{(Quarte)}
\end{equation}
Der Tetraeder enthält beide fundamentalen harmonischen Intervalle:
\begin{equation}
	\frac{3}{2} \times \frac{4}{3} = 2 \quad \text{(Oktave)}
\end{equation}

\section{Schritt 3: Eindeutigkeit von \texorpdfstring{$\xi$}{xi}}

\subsection{Das e-p-\texorpdfstring{$\mu$}{mu}-System als Eindeutigkeitsbeweis}

Aus den drei fundamentalen Massenverhältnissen:
\begin{align}
	R_{pe} &= \frac{m_p}{m_e} = 1836{,}15267343 \\
	R_{\mu e} &= \frac{m_{\mu}}{m_e} = 206{,}7682830 \\
	R_{p\mu} &= \frac{m_p}{m_{\mu}} = 8{,}880
\end{align}
Mit der Massenskalierung $m_p / m_e = 245 \times (4/3)^\kappa$ folgt:
\begin{equation}
	\kappa = \frac{\ln(R_{pe}/245)}{\ln(4/3)} = \frac{\ln(1836{,}15/245)}{\ln(1{,}3333)} \approx 7{,}000
\end{equation}

\subsection{\texorpdfstring{$\kappa = 7$}{kappa = 7} ist die einzige konsistente Lösung}

\begin{table}[h]
\centering
\begin{tabular}{lccr}
\toprule
\textbf{Exponent $\kappa$} & \textbf{$R_{pe}$ Vorhersage} & \textbf{Konsistenz} & \textbf{Fehler} \\
\midrule
$\kappa = 6$ & $245 \times (4/3)^6 = 1376{,}6$ & \texttimes & $25{,}0\%$ \\
$\kappa = 7$ & $245 \times (4/3)^7 = 1835{,}4$ & \checkmark & $0{,}04\%$ \\
$\kappa = 8$ & $245 \times (4/3)^8 = 2447{,}2$ & \texttimes & $33{,}3\%$ \\
\bottomrule
\end{tabular}
\caption{$\kappa = 7$ ist die einzige konsistente Lösung}
\end{table}

\begin{keyresult}[\texorpdfstring{$\xi$}{xi} als emergente Konstante]
\begin{equation}
	\boxed{\xipar = \frac{4}{30000} = \frac{4}{3} \times 10^{-4} = 1{,}333 \times 10^{-4}}
\end{equation}
$\xi$ ist kein Postulat — es ist die einzige selbstkonsistente Lösung des e-p-$\mu$-Systems.
\end{keyresult}

\section{Schritt 4: Gravitationskonstante \texorpdfstring{$G$}{G} aus \texorpdfstring{$\xi$}{xi}}

\subsection{Fundamentale T0-Gravitationsbeziehung}

Der T0-Lagrangian impliziert die fundamentale geometrische Beziehung:
\begin{equation}
	\xipar = 2\sqrt{G \cdot m_{\text{char}}}
	\quad \Rightarrow \quad
	G = \frac{\xipar^2}{4\, m_{\text{char}}} \quad \text{(T0-natürliche Einheiten)}
\end{equation}

Der vollständige SI-Wert erfordert explizite Dimensions- und Einheitenumrechnungsfaktoren (Dokument 012):
\begin{equation}
	\boxed{G_{\text{SI}} = \frac{\xipar^2}{4\, m_e} \times C_{\text{dim}} \times C_{\text{conv}} \times \Kfrak}
	\label{eq:G_from_xi}
\end{equation}
wobei $C_{\text{dim}}$ die Dimensionskorrektur beim Übergang von T0-natürlichen zu physikalischen Einheiten, $C_{\text{conv}}$ den SI-Umrechnungsfaktor und $\Kfrak$ die fraktale Korrektur bezeichnet. Numerisch ergibt sich:
\begin{equation}
	G_{\text{SI}} \approx 6{,}674 \times 10^{-11}\,\text{N\,m}^2/\text{kg}^2
\end{equation}
Abweichung vom CODATA-2018-Wert: $< 0{,}01\%$. Die vollständige Herleitung aller Faktoren findet sich in Dokument 012.

\section{Schritt 5: Minimale Längenskala \texorpdfstring{$L_0$}{L0}}

\subsection{Planck-Länge aus $G$ und \texorpdfstring{$\xi$}{xi}}

\begin{equation}
	\lP = \sqrt{\frac{\hbar G}{c^3}} = 1{,}616 \times 10^{-35}\,\text{m}
\end{equation}
Da $G$ aus $\xi$ abgeleitet ist, folgt $\lP$ vollständig aus $\xi$.

\subsection{\texorpdfstring{$L_0$}{L0} als Sub-Planck-Skala}

Die Zeit-Energie-Dualität $r_0(E) \leftrightarrow E$ in natürlichen Einheiten verlangt eine minimale Längenskala. Der Kopplungsterm $g_T^\ell = \xipar \cdot m_\ell$ skaliert mit $\xipar$ — unterhalb von $L_0$ verliert das Konzept „Abstand" seine physikalische Bedeutung, weil die Zeitfeldfluktuationen $\Delta m$ die Raumzeitstruktur selbst auflösen.

\begin{keyresult}[Minimale Längenskala]
\begin{equation}
	\boxed{L_0 = \xipar \cdot \lP = \frac{4}{30000} \times 1{,}616 \times 10^{-35}\,\text{m} = 2{,}155 \times 10^{-39}\,\text{m}}
	\label{eq:L0}
\end{equation}
\end{keyresult}

\subsection{Physikalische Bedeutung von \texorpdfstring{$L_0$}{L0}}

$L_0$ markiert die Skala, wo T0-geometrische Effekte dominant werden. Unterhalb von $L_0$:
\begin{itemize}
	\item Sind alle Vakuumfluktuationen vollständig wirksam (Casimir-Effekt).
	\item Verliert die störungstheoretische Entwicklung ihre Gültigkeit.
	\item Wird die fraktale Korrektur $\Kfrak$ strukturell notwendig.
\end{itemize}

\section{Zusammenfassung: Vollständige Ableitungskette}

\begin{keyresult}[Ableitungskette ohne Postulate]
\begin{align}
	&\text{Lagrangian-Struktur} \; (\tilde{T} \cdot m = 1) \notag \\
	&\quad\Downarrow \notag \\
	&\text{Selbstkonsistenz verlangt} \; g_T = \xipar \cdot m \notag \\
	&\quad\Downarrow \notag \\
	&10^{-4}\text{-Faktor aus 4D Raumzeit},\quad 4/3\text{-Faktor aus Tetraeder-Geometrie} \notag \\
	&\quad\Downarrow \notag \\
	&\kappa = 7 \text{ einzige Lösung des e-p-}\mu\text{-Systems} \notag \\
	&\quad\Downarrow \notag \\
	&\boxed{\xipar = 4/30000} \quad \text{(kein Postulat — emergente Konstante)} \notag \\
	&\quad\Downarrow \notag \\
	&G = \xipar^2 / (4\, m_e) \quad\Rightarrow\quad \lP = \sqrt{\hbar G/c^3} \notag \\
	&\quad\Downarrow \notag \\
	&\boxed{L_0 = \xipar \cdot \lP = 2{,}155 \times 10^{-39}\,\text{m}}
\end{align}
\end{keyresult}

\section{Verbindung zum SI-System}

\subsection{Die unvermeidliche Zirkularität}

Alle Ergebnisse in diesem Dokument — und in den zugehörigen Dokumenten 181 und 182 — sind zunächst \textbf{verhältnisbasiert}: $\xipar$ ist dimensionslos, $L_0/\lP = \xipar$ ist ein reines Verhältnis, $G = \xipar^2/(4m_e)$ gilt in natürlichen Einheiten. Sobald man aber absolute SI-Zahlenwerte angeben will — also $L_0 = 2{,}155 \times 10^{-39}\,\text{m}$ statt nur $L_0 = \xipar \cdot \lP$ — braucht man mindestens \textbf{eine} externe Verbindung zum SI-System.

\begin{important}[Scheinbare Zirkularität ist unvermeidlich]
Es ist nicht möglich, aus einem rein dimensionslosen Parameter wie $\xipar$ absolute SI-Werte herzuleiten, ohne irgendwo eine Verbindung zur SI-Einheitendefinition herzustellen. Diese Verbindung ist keine Schwäche der Theorie — sie ist eine epistemologische Notwendigkeit jeder physikalischen Theorie.
\end{important}

\subsection{Mögliche Einstiegspunkte ins SI-System}

Die Verbindung kann über verschiedene Größen erfolgen, die alle äquivalent sind:

\begin{itemize}
	\item \textbf{Über $\alpha$}: Die Feinstrukturkonstante $\alpha = \xipar \cdot E_0^2 / (1\,\text{MeV})^2$ ist dimensionslos und experimentell sehr genau bekannt. Sie liefert über $E_0 = \sqrt{m_e \cdot m_\mu}$ die Verbindung zu den Leptonmassen in MeV.
	\item \textbf{Über $m_{\text{char}}$ bzw.\ $m_e$}: Die Elektronmasse $m_e = 0{,}511\,\text{MeV}$ ist experimentell definiert und verknüpft $\xipar$ mit SI-Einheiten über $G = \xipar^2/(4m_e)$.
	\item \textbf{Über $\lP$}: Die Planck-Länge $\lP = \sqrt{\hbar G/c^3}$ ist experimentell zugänglich und liefert direkt $L_0 = \xipar \cdot \lP$ in Metern.
\end{itemize}

\subsection{Welcher Einstieg ist fundamental?}

In T0/FFGFT ist der natürlichste Einstieg über $\alpha$, weil:
\begin{itemize}
	\item $\alpha$ selbst aus $\xipar$ abgeleitet wird: $\alpha = \xipar \cdot E_0^2 / (1\,\text{MeV})^2$
	\item $\alpha$ dimensionslos ist — der Übergang zu SI-Einheiten erfolgt erst durch die Wahl der Einheitenskala (MeV)
	\item Die SI-Reform 2019 hat genau diese Kalibration implementiert, die mit der T0-Geometrie konsistent ist (Dokument 013)
\end{itemize}

\begin{keyresult}[Epistemologische Einordnung]
Die Ableitungskette in diesem Dokument ist intern konsistent und postulat-frei. Die einzige externe Eingabe ist die Wahl einer SI-Einheitenskala — zum Beispiel über $m_e$ in MeV oder $\lP$ in Metern. Diese Eingabe ist keine Zirkularität der Theorie, sondern die notwendige Kalibrierung jeder physikalischen Theorie an die Messwelt. Ausführlicher behandelt in Dokument 056 (Universale Ableitung) und Dokument 101 (Zirkularität der Konstanten).
\end{keyresult}


